\section{Spesifikasi Token Subset-C}
\label{sec:ch05-subsetc}

\subsection{Kategori Token}

Subset-C memiliki tujuh kategori token: kata kunci, identifier,
literal, operator, tanda baca, akhir masukan, dan kesalahan. Daftar
lengkap terdapat pada \file{code/subsetc/spec/tokens.md}. Tabel
berikut merangkum kategori dan contohnya.

\begin{table}[H]
\centering
\caption{Ringkasan kategori token Subset-C}
\label{tab:ch05-subsetc}
\begin{tabularx}{\textwidth}{@{}l l X@{}}
\toprule
\textbf{Kategori} & \textbf{Contoh token} & \textbf{Contoh lexeme} \\
\midrule
Kata kunci & \code{TOK\_INT}, \code{TOK\_IF}, \code{TOK\_WHILE} &
\code{int}, \code{if}, \code{while} \\
Identifier & \code{TOK\_IDENT} & \code{hasil}, \code{\_x1} \\
Literal & \code{TOK\_INT\_LIT}, \code{TOK\_STRING\_LIT} &
\code{42}, \code{"halo"} \\
Operator & \code{TOK\_LE}, \code{TOK\_AND} & \code{<=}, \code{\&\&} \\
Tanda baca & \code{TOK\_LPAREN}, \code{TOK\_SEMI} & \code{(}, \code{;} \\
Akhir & \code{TOK\_EOF} & --- \\
Kesalahan & \code{TOK\_ERROR} & \code{@} \\
\bottomrule
\end{tabularx}
\end{table}

\subsection{Header Token}

Header \file{token.h} mendefinisikan enum \code{TokenType} dan
struktur \code{Token}. Enum memberi nama simbolik pada setiap jenis
token, sedangkan struktur menyimpan lexeme, lokasi, dan nilai.
Bagian awal header ditunjukkan pada \lstref{lst:ch05-token-h}.

\includecode[
  language=C,
  firstline=8,
  lastline=32,
  firstnumber=8,
  caption={Bagian awal \file{token.h}: enum jenis token},
  label={lst:ch05-token-h}
]{../../code/subsetc/include/token.h}

Setiap token memiliki nama yang dapat dicetak melalui
\code{token\_type\_name}. Fungsi ini berguna untuk keluaran tes dan
pesan diagnostik. Struktur \code{Token} menyimpan \code{lexeme}
sebagai pointer ke buffer yang dimiliki lexer; pemanggil tidak boleh
membebaskannya.

\subsection{Program Uji}

Program \file{lex\_demo.sc} dirancang untuk menguji seluruh kategori
token, termasuk komentar, literal, dan operator dua karakter. Potongan
awalnya ditunjukkan pada \lstref{lst:ch05-lex-demo}.

\includecode[
  language=C,
  firstline=1,
  lastline=12,
  firstnumber=1,
  caption={Potongan program uji \file{lex\_demo.sc}},
  label={lst:ch05-lex-demo}
]{../../code/subsetc/examples/lex_demo.sc}

Program ini akan dipakai pada Bab VII untuk membandingkan keluaran
lexer tangan dan lexer Flex. Hasil yang diharapkan adalah aliran token
yang identik, kecuali perbedaan lokasi yang disebabkan oleh penanganan
baris baru.

\subsection{Latihan Spesifikasi}

Sebelum menulis kode, lengkapi spesifikasi token dengan:
\begin{itemize}[leftmargin=*]
  \item tabel pattern untuk setiap token;
  \item daftar kata kunci;
  \item aturan escape untuk char dan string;
  \item perilaku untuk karakter ilegal;
  \item minimal 30 kasus uji, termasuk kasus batas.
\end{itemize}

Spesifikasi yang lengkap akan mempermudah implementasi dan pengujian.
Sebaliknya, spesifikasi yang ambigu biasanya berujung pada lexer yang
berperilaku berbeda dari yang diharapkan.
